\volumeopening{About This Volume}

\section*{Why Study the Properties of Matter and Heat?}

Volume~I treated bodies as particles or as perfectly rigid objects, and that idealisation carried it a
long way. Real matter is neither. A steel cable stretches under load, a ship floats because water pushes
back, a railway line lengthens on a summer afternoon, and a gas heats up when it is compressed. This
volume asks how bulk matter behaves, and it answers from two directions. First, how do solids and fluids
respond to \emph{force}: what makes a cable strong, a ship float, a wing lift? Then, how does matter
respond to \emph{heat}: what makes a metal expand, water boil, an engine turn?

The second question has shaped history more than almost any other in physics. The steam engine, a
machine that turns heat into motion, powered the Industrial Revolution, and the effort to understand what
limited it created thermodynamics. Every engine, refrigerator, power station and jet turbine running today
is an application of the handful of laws in these pages, and the question ``how much useful work can be
obtained from a given amount of heat?'' still sets the efficiency of the electricity grid.

The same laws also say something about time. Thermodynamics is where physics first admitted that some
processes run one way only: heat flows from hot to cold and never back unaided, a broken cup does not
reassemble, and scent spreads through a room but never gathers back into the bottle. This one-way
character, captured by the second law, explains why no engine can turn heat entirely into work and why
perpetual-motion machines are impossible.

The ideas reach well beyond physics. Engineers use them to size beams, cables, engines and power plants;
materials scientists, to explain why bridges need expansion joints and why frozen pipes burst; biologists,
to understand the regulation of body temperature; earth scientists, to model ocean currents and the
greenhouse effect; chemists, to describe the energetics of reactions and the behaviour of gases; and
astrophysicists, to explain how a star balances gravity against the pressure of its hot interior.

\section*{How This Volume Is Built}

The volume has five chapters in three parts, and each part answers one question.

\textbf{Part A, Mechanical Properties of Matter,} asks how bulk matter responds to force.
\chref{ch:solids} studies how solids stretch, compress, shear and twist, why the response is described by
stress and strain rather than force and extension, and when a solid stops springing back and breaks.
\chref{ch:fluids} turns to fluids, which cannot resist a shear at all and so flow: pressure, buoyancy,
the continuity and Bernoulli equations, viscosity and surface tension.

\textbf{Part B, Heat and Temperature,} asks what heat does to ordinary matter. \chref{ch:thermal} defines
temperature, then follows heat as it raises temperature, expands solids and liquids, melts and boils
substances, and flows from hot to cold by conduction, convection and radiation.

\textbf{Part C, Thermodynamics and the Molecular Picture,} asks what rules govern heat and work, and why.
\chref{ch:thermo} states the first law (energy is conserved, with heat counted as energy in transit) and
the second (heat has a preferred direction), and derives the efficiency of the best possible engine.
\chref{ch:kinetic} then looks inside the gas and shows that pressure, temperature and the gas laws all
emerge from molecules in motion obeying Newton's laws.

Each idea is introduced by why it is worth having before how it works, and each result is derived or, where
a derivation is beyond the level of the book, flagged as stated without proof. By the end you will be able
to explain, from first principles, why a thick rope holds more than a thin one, why a ship of steel floats,
why an aeroplane wing lifts, why a railway line buckles in summer, why a pressure cooker cooks faster, why
no engine can be perfectly efficient, and why a gas warms when it is compressed.

\section*{What It Builds On}

The volume is the mechanics of Volume~I applied to matter in bulk, and it leans on that volume
throughout. Hooke's law and the spring constant (\chref{ch:connected}) become Young's modulus;
the work done by a varying force and elastic potential energy (\chref{ch:energy}) become the energy
stored in a stretched wire and the work done by a gas; Newton's second law and the momentum of colliding
particles (\chref{ch:newton} and \chref{ch:collisions}) give the pressure of a gas; and the moment of
inertia (\chref{ch:rotation}) reappears in the stiffness of a beam. When a quantity from Volume~I returns it
keeps its Volume~I symbol: kinetic energy $K$, length $\ell$, spring constant $k$.

The mathematics is collected in \chref{ch:toolkit}, at the front of the book. Besides
derivatives, integrals and the binomial approximation $(1+x)^n\approx1+nx$ (Section~\ref{sec:approx}),
which this volume uses on almost every page, two of its sections were written with this volume in mind:
exponentials, logarithms and first-order rate equations (Section~\ref{sec:explog}), needed for Newton's law
of cooling and for isothermal work, and averages over a probability distribution with the two Gaussian
integrals (Section~\ref{sec:averages}), needed for the speeds of gas molecules. Read them when a derivation
first calls on them, or now if they are unfamiliar.

\clearpage
\section*{Notation Used in This Volume}

Thermal physics is hard on the alphabet, and this volume keeps to one rule: within the volume, one symbol
means one quantity. The table lists the symbols and where each first does its work; the paragraphs after
it single out the collisions that most often cause examination errors.

\begin{center}\small
\begin{tabular}{@{}lp{0.38\textwidth}lp{0.27\textwidth}@{}}
\toprule
Symbol & Meaning (unit) & Where & Note\\
\midrule
$\sigma$ & stress, force per area (\si{Pa}) & \chref{ch:solids} & not Stefan's constant\\
$\epsilon$ & strain, fractional deformation & \chref{ch:solids} & dimensionless\\
$Y,\ G,\ B$ & Young's, shear and bulk moduli (\si{Pa}) & \chref{ch:solids} & $B$, to keep $K$ for conductivity\\
$\nu$ & Poisson's ratio & \chref{ch:solids} & dimensionless\\
$\ell,\ \Delta\ell$ & length of a wire or rod; its extension (\si{m}) & \chref{ch:solids} & as in Volume~I\\
$k$ & spring constant (\si{N/m}) & \chref{ch:solids} & from Volume~I\\
$I_g$ & second moment of area of a beam's section (\si{m^4}) & \chref{ch:solids} & not a moment of inertia\\
$P,\ P_0;\ \rho$ & pressure, atmospheric pressure \SI{1.013e5}{Pa}; density (\si{kg/m^3}) & \chref{ch:fluids} & \\
$\eta$ & coefficient of viscosity (\si{Pa.s}) & \chref{ch:fluids} & see the cautions\\
$S$ & surface tension (\si{N/m}, equivalently \si{J/m^2}) & \chref{ch:fluids} & never $T$, which is temperature\\
\midrule
$T$ & absolute temperature (\si{K}) & \chref{ch:thermal} & only temperature\\
$t_C,\ t$ & Celsius temperature (\si{\celsius}); time (\si{s}) & \chref{ch:thermal} & \\
$\alpha_\ell,\ \alpha_A,\ \alpha_V$ & linear, area and volume expansion coefficients (\si{K^{-1}}) & \chref{ch:thermal} & never $\beta$ or $\gamma$\\
$Q$ & heat supplied \emph{to} the system (\si{J}) & \chref{ch:thermal} & energy in transit\\
$s$ & specific heat capacity (\si{J.kg^{-1}.K^{-1}}) & \chref{ch:thermal} & per kilogram\\
$L_f,\ L_v$ & latent heats of fusion, vaporisation (\si{J/kg}) & \chref{ch:thermal} & \\
$K$ & thermal conductivity (\si{W.m^{-1}.K^{-1}}) & \chref{ch:thermal} & capital\\
$H$ & heat current $dQ/dt$ (\si{W}) & \chref{ch:thermal} & not a height\\
$R_{\text{th}}$ & thermal resistance (\si{K/W}) & \chref{ch:thermal} & \\
$\theta$ & an angle (shear strain, angle of contact) in Part~A; temperature excess $T-T_a$ in \chref{ch:thermal} & \chref{ch:solids} & see the cautions\\
$\kappa$ & cooling constant in Newton's law (\si{s^{-1}}) & \chref{ch:thermal} & not plain $k$\\
$\sigma_{\!s},\ e$ & Stefan's constant; emissivity & \chref{ch:thermal} & \\
$n,\ R,\ N_A$ & amount of substance (\si{mol}); gas constant; Avogadro's number & \chref{ch:thermal} & $n$ always moles\\
\midrule
$W$ & work done \emph{by} the system (\si{J}) & \chref{ch:thermo} & sign convention\\
$U$ & internal energy (\si{J}) & \chref{ch:thermo} & \\
$C,\ C_p,\ C_v$ & molar heat capacity ($C=Ms$, met in \chref{ch:thermal}); its values at constant pressure and volume (\si{J.mol^{-1}.K^{-1}}) & \chref{ch:thermo} & per mole\\
$\gamma$ & ratio $C_p/C_v$ & \chref{ch:thermo} & only this\\
$x,\ C_x$ & polytropic index ($PV^{x}=$ constant); its molar heat capacity & \chref{ch:thermo} & not $n$\\
$T_h,\ T_c$ & temperatures of the hot and cold reservoirs (\si{K}) & \chref{ch:thermo} & \\
$Q_h,\ Q_c$ & heat taken from the hot, given to the cold reservoir (\si{J}) & \chref{ch:thermo} & magnitudes\\
$\eta$;\ COP & efficiency of an engine; coefficient of performance of a refrigerator & \chref{ch:thermo} & see the cautions\\
$S$ & entropy (\si{J/K}) & \chref{ch:thermo} & see the cautions\\
\midrule
$k_B$ & Boltzmann's constant (\si{J/K}) & \chref{ch:kinetic} & always subscripted\\
$N$ & number of molecules, $nN_A$ & \chref{ch:kinetic} & \\
$n_V$ & number density $N/V$ (\si{m^{-3}}) & \chref{ch:kinetic} & \\
$m,\ M$ & mass of one molecule; molar mass & \chref{ch:kinetic} & \\
$\phi(v)$ & distribution of molecular speeds (\si{s/m}) & \chref{ch:kinetic} & \\
$\overline{K}$ & mean translational kinetic energy of a molecule (\si{J}) & \chref{ch:kinetic} & \\
$v_p,\ \bar v,\ v_{\text{rms}}$ & most probable, mean and root-mean-square speeds (\si{m/s}) & \chref{ch:kinetic} & \\
$f$ & number of degrees of freedom & \chref{ch:kinetic} & \\
$\lambda$ & mean free path (\si{m}) & \chref{ch:kinetic} & not a wavelength\\
$d$ & molecular diameter (\si{m}) & \chref{ch:kinetic} & \\
$z$ & collision frequency of one molecule (\si{s^{-1}}) & \chref{ch:kinetic} & \\
\bottomrule
\end{tabular}
\end{center}

\begin{trapbox}{the symbol collisions that cause examination errors}
\textbf{(i) $\gamma$.} Many books use $\gamma$ for the volume expansion coefficient in thermal expansion
\emph{and} for $C_p/C_v$ in thermodynamics. This book does not: the expansion coefficients are
$\alpha_\ell$, $\alpha_A$, $\alpha_V$ (with $\alpha_A=2\alpha_\ell$ and $\alpha_V=3\alpha_\ell$, derived in
\chref{ch:thermal}), and $\gamma$ always means $C_p/C_v$.\par\smallskip
\textbf{(ii) $s$ and $C$.} Specific heat $s$ is per \emph{kilogram}, molar heat capacity $C$ per
\emph{mole}; they differ by the molar mass, $C=Ms$. Mixing them up shifts an answer by a factor of
$10$ to $1000$.\par\smallskip
\textbf{(iii) $K$, $B$, $k$ and $k_B$.} Capital $K$ is thermal conductivity, which is why the bulk modulus
is written $B$. Boltzmann's constant is always $k_B$. Plain $k$ keeps its Volume~I meaning, the spring
constant, and appears only when a stretched wire is treated as a spring in \chref{ch:solids} and a molecular bond as a spring in \chref{ch:kinetic}; the cooling
constant of Newton's law of cooling is written $\kappa$ (the generic rate constant $k$ of the
Foundations' equation $dy/dt=-ky$).
\par\smallskip
\textbf{(iv) $n$ and $x$; $f$ and $\phi$.} The letter $n$ always counts moles, so the index of a polytropic
process $PV^{x}=$ constant is written $x$. And $f$ counts degrees of freedom, so the distribution of
molecular speeds is written $\phi(v)$.\par\smallskip
\textbf{(v) $\sigma$.} Plain $\sigma$ is a mechanical stress; Stefan's radiation constant carries a
subscript, $\sigma_{\!s}$, so the two never clash.\par\smallskip
\textbf{(vi) $\eta$, $S$ and $\theta$.} The viscosity of a fluid and the efficiency of an engine are both traditionally $\eta$; surface tension (in this book, following NCERT) and entropy are both $S$; and $\theta$ is an angle in Part~A but the temperature excess of Newton's law of cooling in \chref{ch:thermal}. In each pair the two meanings never appear in the same chapter.\par\smallskip
\textbf{(vii) $U$.} $U$, the potential energy of Volume~I, here means the internal energy of a system, which includes the potential energy of
its molecules but also their kinetic energy.
\end{trapbox}
